% a06_BALtorus_ct3: Step 3(E), third part: the finite part dm^2_fin (Ba3 (1.26)) in the form in which it is consumed,
% and the size of the non-renormalising (1.7) vertex at every step (the vertex tolerance).

\subsubsection*{The finite part $\dm_{\mathrm{fin}}$}

By \cref{lem:a06e:MC}(a) with $H=(G_K,G^v_K)$ and $D=(R_K,R^v_K)$, $\dm_{\mathrm{fin}}=\dm[C_0,C]-\dm[G_K,G^v_K]$ is a finite sum of terms in which every cluster is a renormalised family at step $K$ with lines $G_K$ and every cut line is an $R_K$ line. The order-$2$ part is \cref{thm:a06e:DM}(e). The order-$4$ part is bounded in \cref{thm:a06e:DMf} below, up to a polylogarithm; the proof is Proposition~1 at step $K$ per cluster, with the clusters glued along the $R_K$ lines by an infrared power count.

\begin{lemma}[the remainder in scale-$K$ units]\label{lem:a06e:RKp}
Let $\alpha_0\in(0,1)$ and $t\defeq t_K=l_K$. \textup{(a)} For $\sigma,\sigma'\in\{0,1\}$, $|\partial^\sigma\partial'^{\sigma'}R_K(x,x')|\le C_R\max(l_K,|y|)^{-1-\sigma-\sigma'}e^{-c_*|y|}$ (\cref{lem:a06e:RK}), and the H\"older quotients of exponent $\alpha_0$ (in units of $l_K$, over separations $\le2l_K$) of $\partial^\sigma\partial'^{\sigma'}R_K$, $\sigma+\sigma'\le2$, in either variable obey the same bound with $\max(l_K,|y|)^{-1-\sigma-\sigma'}$ (so the mixed quotient, $\sigma=\sigma'=1$, carries $\max(l_K,|y|)^{-3}e^{-c_*|y|}$). \textup{(b)} In scale-$K$ units, $\hat R\defeq l_KR_K(l_K\cdot,l_K\cdot)$: for localisations $h,h'$ at boxes of side $\le2$ at distance $r$, $\norm{(h\otimes h')\partial^\sigma\partial'^{\sigma'}\hat R}_{0,\alpha_0}\le C_R\max(1,r)^{-1-\sigma-\sigma'}e^{-c_*tr}$, and the two-variable norm $\norm{(h\otimes h')\hat R}_{1,\alpha_0}$ is $\le C_R\max(1,r)^{-1}e^{-c_*tr}$. \textup{(c)} The same for $R^v_K$. The constants are free of $\eps$, $K$ and the torus.
\end{lemma}

\begin{proof}
(a) The sup and first-difference bounds are \cref{lem:a06e:RK}. For the H\"older quotients use $R_K=G_K(a_Kl_K^{-2}P_K)C_0$ \cite[(1.25)]{Ba3}: with the block average, $R_K(x,x')=a_Kl_K^{-5}\sum_B(G_K\mathbf 1_B)(x)u_B(x')$, $u_B(x')\defeq\sum_{z\in B}\eps^3C_0(z,x')$, the sum over the blocks of side $l_K$. Since $G_K\mathbf 1=(m^2+a_Kl_K^{-2})^{-1}\mathbf 1$ on the torus, $\sum_B\partial_x(G_K\mathbf 1_B)(x)=0$ and $\partial_xR_K(x,x')=a_Kl_K^{-5}\sum_B\partial_x(G_K\mathbf 1_B)(x)[u_B(x')-u_{B_x}(x')]$, $B_x$ the block of $x$. By \cite[Lemma~2.4, (2.35)--(2.36)]{Ba4} at scale $K$ in the form of \cref{thm:a06:A5} ($\mathbf 1_B=Q_K^*\delta_{\hat y}$ in scale-$K$ units; the torus is a periodic box that is a union of the blocks of side $l_K$, so no depth condition arises and the constants are free of the torus) rescaled to $\eps$-units, $|G_K\mathbf 1_B|(x)\le c_0l_K^2e^{-\delta_0\dist(x,B)/l_K}$, $|\partial_xG_K\mathbf 1_B|\le c_0l_Ke^{-\delta_0\dist(x,B)/l_K}$, and the H\"older quotient of $\partial_xG_K\mathbf 1_B$ (in units of $l_K$) obeys the same bound with $c_1(\alpha_0)e^{2\delta_0}$ in place of $c_0$, $c_1(\alpha_0)$ the constant of \cite[(2.36)]{Ba4} (which decays in $\dist(\{\hat x,\hat x'\},\hat y)\ge\dist(\hat x,B)-2$ for separations $\le2l_K$). By \cref{lem:a06e:SL}(c) along a path from $B_x$ to $B$, for $|y|\ge4l_K+2\dist(B,B_x)$: $|u_B(x')-u_{B_x}(x')|\le Cl_K^3(\dist(B,B_x)+l_K)|y|^{-2}e^{-c_*|y|}$, each difference in $x'$ adding $|y|^{-1}$; for the remaining $B$ ($\dist(B,B_x)\ge|y|/4-l_K$) bound the two terms separately by $|u_B|\le Cl_K^3\max(l_K,\dist(B,x'))^{-1}$. The sum over $B$ converges by $e^{-\delta_0\dist/l_K}$, which absorbs the polynomial factors and makes the remaining-$B$ part $\le Ce^{-\delta_0|y|/(8l_K)}$. This gives (a) for the H\"older quotient of $\partial_xR_K$ at $|y|\ge4l_K$; for $|y|<4l_K$ the same sum with $|u_B|\le Cl_K^2$ gives $Cl_K^{-2-\sigma'}$. The $x'$-variable is symmetric ($R_K=C_0a_Kl_K^{-2}P_KG_K$). The mixed quotient (the H\"older quotient of $\partial_x\partial_{x'}R_K$ in either variable) is the same sum with one more difference: \cref{lem:a06e:SL}(c) at $\sigma+\sigma'\le3$ for the differences of $u_B$, and \cite[(2.36)]{Ba4} (in the form of \cref{thm:a06:A5}) for $\partial_xG_K\mathbf 1_B$. (b) is (a) rescaled: a two-point kernel scales by $l_K$, each difference by $l_K$, and $e^{-c_*|y|}=e^{-c_*tr}$; the norm $\norm{\cdot}_{1,\alpha_0}$ is the maximum of the sup, the first differences and their H\"older quotients.
\end{proof}

\begin{lemma}[derivative credit]\label{lem:a06e:DC}
Let a localised family at step $K$ (zero background, torus) have an external leg $l$ at a vertex $v$ of type \cite[(1.8)]{Ba3}, with the difference of $v$ acting on $l$, and let $f_l$ be the external field on $l$. Then \cite[(1.33)]{Ba3} holds with $\norm{hf_l}_{1,\alpha_0}$ replaced by $C_{\mathrm{DC}}\norm{\hat h\,\partial f_l}_{0,\alpha_0}$, $\hat h$ a localisation equal to $1$ on the $1$-neighbourhood of $\supp h$, $C_{\mathrm{DC}}$ depending on $(d,\alpha_0)$ and the localisation family only; the same in each variable of a two-variable external field.
\end{lemma}

\begin{proof}
Vertices are localised by weights on their summation variables (unit-cube indicators for \cite[(1.6)--(1.7)]{Ba3}, smooth partitions of unity on the vector legs for the others, \cite[p.~420]{Ba3}), so $f_l$ enters only through $(\partial f_l)(b)$ at the bonds of the localisation of $v$, weighted by the partition, and the value is unchanged under $f_l\to f_l-f_l(z_0)$, $z_0\in\supp h$ (at zero background constants have zero difference). Apply \cite[(1.33)]{Ba3} to $f_l-f_l(z_0)$: on $\supp h$, $|f_l(z)-f_l(z_0)|\le3d\sup|\partial f_l|$ along an axis path, and $\norm{h(f_l-f_l(z_0))}_{1,\alpha_0}\le C(\sup_{\supp\hat h}|f_l-f_l(z_0)|+\norm{\hat h\,\partial f_l}_{0,\alpha_0})$.
\end{proof}

\begin{lemma}[infrared power counting for the gluing]\label{lem:a06e:IRC}
Let $Q$ be a finite connected multigraph (self-edges allowed) with a root node $c_0$; each non-root node $c$ has a position $z_c\in\Z^3$, $z_{c_0}\defeq0$ (or positions on a torus, with the torus distance in the weights, torus balls in the count and the torus triangle inequality; every step below holds verbatim). A non-self edge $e$ between $c,c'$ has weight $w_e(z_c-z_{c'})$ with $0\le w_e(z)\le\max(1,|z|)^{-a_e}e^{-\mu|z|}$, $a_e\ge0$, $0<\mu\le\frac12$; self-edges have weight $\le1$. For a partition $P$ of the nodes into groups each inducing a connected subgraph put $\mathrm{IR}(P)\defeq3(|P|-1)-\sum_{e\text{ between groups}}a_e$ and $D\defeq\max_P\mathrm{IR}(P)\ge0$. Then $\sum_{(z_c)}\prod_ew_e\le C_Q(1+\log_2\mu^{-1})^{|E|}\mu^{-D}$, $C_Q$ depending on $Q$ and $(a_e)$ only.
\end{lemma}

\begin{proof}
For a non-self edge put $j_e\defeq0$ if $|z_e|\le1$ and $j_e\defeq j$ if $2^{j-1}<|z_e|\le2^j$; $J\defeq\lceil\log_2\mu^{-1}\rceil$, $\hat j_e\defeq\min(j_e,J)$. Kruskal's rule with the edges in non-decreasing $j_e$ gives a spanning tree $\tau$; the configuration is fixed by the tree-edge vectors, so the number of configurations with a given scale vector is $\le\prod_{e\in\tau}C2^{3j_e}$. Weights: $w_e\le2^{a_e}2^{-a_ej_e}e^{-\mu2^{j_e-1}[j_e\ge1]}$. For $e\in\tau$ with $j_e>J$: $2^{(3-a_e)j_e}e^{-\mu2^{j_e-1}}\le2^{(3-a_e)\hat j_e}g(j_e-J)$ with $g(n)\defeq2^{3n}\exp(-2^{n-2})$ summable; for $e\notin\tau$, $2^{-a_ej_e}\le2^{-a_e\hat j_e}$. Telescoping: $\sum_{e\in\tau}3\hat j_e-\sum_ea_e\hat j_e=\sum_{i=1}^J[3\#\{e\in\tau:j_e\ge i\}-\sum_{e:j_e\ge i}a_e]$; for level $i$ let $P_i$ be the partition into the components of the edges with $j_e<i$; by Kruskal's rule $\#\{e\in\tau:j_e\ge i\}=|P_i|-1$, every edge between groups of $P_i$ has $j_e\ge i$, and the other edges with $j_e\ge i$ only lower the bracket; so the bracket is $\le\mathrm{IR}(P_i)\le D$ and the total is $\le JD$, i.e.\ $2^{JD}\le(2/\mu)^D$. A rejected edge $e$ closes a cycle of tree edges $f$ with $j_f\le j_e$ (Kruskal's rule), and the triangle inequality along that path gives $2^{j_e-1}<|z_e|\le|Q|2^{\max_fj_f}$, so $\max_fj_f\le j_e\le\max_fj_f+\log_2|Q|+1$; hence each tree scale vector extends to at most $(\log_2|Q|+2)^{|E|-|\tau|}$ full scale vectors, there are $\le(J+1)^{|\tau|}$ truncated tree scale vectors, and the excess sums converge; the stated exponent $|E|\ge|\tau|$ is an upper bound.
\end{proof}

\begin{lemma}[degrees]\label{lem:a06e:ID}
For a connected subgraph $g$ of a graph of Step~3(E) (vertex set and all lines among them) let $E(g)$ be the number of its legs not on its lines, $x(g)$ the number of its $V^8$ vertices whose difference acts on such a leg, and $\omega(g)\defeq3\mathcal L(g)-2\sum_{l\in g}(1-d_l/2)-\sum_{v\in g}s_v$ with $s(V_\delta)\defeq-1$. \textup{(a)} $\omega(g)=3-E(g)/2-\mathrm{ord}(g)/2-x(g)$, an integer. \textup{(b)} (Gluing identity.) Let $\Gamma$ be a graph of $\dm_{(4)}$ at zero external momentum, rooted at $x$, $\Lambda$ a set of its lines, the clusters the components of $\Gamma\setminus\Lambda$; give an $R_K$ line $e$ the exponent $a_e\defeq1+d_e$, $d_e\in\{0,1,2\}$ the number of \cite[(1.8)]{Ba3} differences acting on its ends. For a partition $P$ of the clusters into groups connected through $\Lambda$, with $g_1,\dots,g_n$ the subgraphs induced by the groups, $\mathrm{IR}(P)=\sum_i\omega(g_i)-\omega(\Gamma)$. \textup{(c)} (Subgraph degrees.) If $\Gamma$ is a graph of $\dm_{(4)}$ (two-point, order $4$, one-particle irreducible, no difference on its external legs), then $\omega(\Gamma)=0$, every connected subgraph $g\ne\Gamma$ has $\omega(g)\le1$, with $\omega(g)=1$ only if $g$ is blob-shaped ($\mathrm{ord}(g)=2$, $E(g)=2$, $x(g)=0$) and $\omega(g)\le0$ otherwise, and no two vertex-disjoint connected subgraphs of $\Gamma$ are both blob-shaped.
\end{lemma}

\begin{proof}
(a) With $\mathcal L=I-V+1$ and $2I=\sum_vh_v-E$: $\omega=\sum_v(h_v/2-3-s_v+d_{v,g})-E/2+3$, $d_{v,g}\in\{0,1\}$ indicating a $V^8$ difference acting on a line of $g$; the bracket is $-\mathrm{ord}(v)/2$ for $V_\lambda$, $V^{10}_n$, $V_\delta$, and $-n/2-(1-d_{v,g})$ for $V^8_n$ (\cref{lem:a06e:M3}(d)). Parity: $h_v\equiv\mathrm{ord}(v)\pmod2$ for every type, so $E(g)\equiv\mathrm{ord}(g)$ and $\omega(g)\in\Z$. (b) Lines between groups are $R_K$ lines; $\mathcal L(\Gamma)=\sum_i\mathcal L(g_i)+(E_{\mathrm{inter}}-n+1)$, so $\omega(\Gamma)=\sum_i\omega(g_i)+3(E_{\mathrm{inter}}-n+1)-\sum_{\mathrm{inter}}(2-d_e)=\sum_i\omega(g_i)-[3(n-1)-\sum_{\mathrm{inter}}(1+d_e)]$. (c) $\omega(\Gamma)=3-1-2-0=0$. For a proper connected $g$: $E(g)\ge2$ ($E=0$ forces $g=\Gamma$; $E=1$ makes the single leg a bridge or disconnects $\Gamma$). The only order-$1$ vertex is $V^8_1$, whose single-vertex subgraph has $E=3$ and $x=1$. So $E+\mathrm{ord}+2x\ge4$ with equality only for $(E,\mathrm{ord},x)=(2,2,0)$: $\omega=3-(E+\mathrm{ord}+2x)/2\le1$, with equality exactly for blob-shaped $g$ and $\omega\le0$ otherwise (integrality). Two vertex-disjoint blob-shaped $g_1,g_2$ carry all of the order, hence all vertices, and each has exactly two external legs; the two external legs of $\Gamma$ are among these four: if both sit on $g_1$, the two legs of $g_2$ must join $g_1$, which has none left; if one on each, the remaining two form a bridge, against one-particle irreducibility; if neither, $\Gamma$ has no external leg.
\end{proof}

\begin{theorem}[the order-$4$ finite part]\label{thm:a06e:DMf}
For every $x$ and every torus aligned with the blocks with sides $\ge1$, $|{\dm}^{(4)}_{\mathrm{fin}}(x)|\le C_{\mathrm{DMf}}\,l_K^{-1}(1+\ln l_K^{-1})^5$, with $C_{\mathrm{DMf}}$ free of $\eps$, $\eps_0$, $K$ and the torus.
\end{theorem}

\begin{proof}
(1) By \cref{lem:a06e:MC}(a), ${\dm}^{(4)}_{\mathrm{fin}}(x)=\sum_\Gamma\sum_{\Lambda\ne\emptyset}(\Gamma\setminus\Lambda)_{\mathrm{ren}}[G_K;R_K\text{ on }\Lambda](x)$, a finite sum (at most $4$ vertices and $5$ lines per $\Gamma$). (2) Each cluster is a renormalised family at step $K$ with lines $G_K$; its \cite[(1.7)]{Ba3} vertices are $\dm_1$ vertices (no order-$4$ counterterm occurs in $\dm_{(4)}$), whose non-renormalising part is $\Delta_{K,K}=O(l_K^{-1})$ (\cref{prop:a06e:DM1}), of the printed size; in scale-$K$ units Proposition~1 gives per localised cluster family $O(1)t^{o_c/2}e^{-\delta'd(\{\square\})}$ times the norms of its external fields, $\sum_co_c=4$. (3) The external fields are the ends of the $R_K$ lines; an end at which a \cite[(1.8)]{Ba3} difference acts is measured by the differenced kernel (\cref{lem:a06e:DC}); by \cref{lem:a06e:RKp}(b) a line $e$ between boxes at distance $r$ contributes $\le C_R\max(1,r)^{-a_e}e^{-c_*tr}$, $a_e=1+d_e$, and a line with both ends in one cluster $\le C_R$; clusters joined by $R_K$ lines are joined as in \cref{lem:a06d:V10}. (4) For each cluster sum its boxes around a root box with half of its tree decay (\cref{lem:a06d:V11}); the distance of two attached boxes differs from that of the roots by at most $d_c+d_{c'}$, and $(1+d_c+d_{c'})^{a_e}e^{c_*t(d_c+d_{c'})}$ is absorbed by the other half ($c_*t\le\delta'/8$). (5) \Cref{lem:a06e:IRC} with $\mu\defeq c_*t$ bounds the sum over the cluster roots by $C(1+\log t^{-1})^{|\Lambda|}t^{-D}$, and by \cref{lem:a06e:ID}(b),(c), $\mathrm{IR}(P)=\sum_i\omega(g_i)-\omega(\Gamma)\le1$ for every partition (at most one blob-shaped group, every other group $\le0$), so $D\le1$. (6) In scale-$K$ units the coefficient is $\dm(x)l_K^2$, so $|{\dm}^{(4)}_{\mathrm{fin}}(x)|l_K^2\le Ct^2t^{-1}(1+\log t^{-1})^5$ with $t=l_K$.
\end{proof}

The logarithm is the count of scale vectors in \cref{lem:a06e:IRC}; the statement without it is not claimed, and nothing below needs it.

\subsubsection*{The non-renormalising vertex at every step}

For $0\le j\le K$ put $c_j(x)\defeq l_j^2(\dm_{\mathrm{fin}}(x)+\dm_{K,j}(x))$, $\dm_{K,K}\defeq0$: the coefficient, in scale-$j$ units, of the part of the step-$j$ vertex \cite[(1.7), (1.29)]{Ba3} that is not absorbed into renormalised pairs ($c_0=\eps^2\dm$, since $\dm_0=0$); the renormalising piece $\dm_j(e_j,\lambda_j,\Omega_1,x)$ of \cite[(1.29)]{Ba3} is of size $l_j^2/\eps$-type by \cref{thm:a06e:CT1} rescaled and is cancelled inside the renormalised classes. Let $q(t)\defeq(1+\ln t^{-1})^5$ and $Q_\beta\defeq\sup_{0<t\le1}t^\beta q(t)=(5/\beta)^5e^{\beta-5}$ for $\beta\in(0,5]$.

\begin{theorem}[the vertex tolerance]\label{thm:a06e:T0}
Let $C_T\defeq2\max(C_D,c_m)+2C_4/\ln L+C_{\mathrm{DMf}}$. For every $0\le j\le K$ and $x$: \textup{(a)} $|c_j(x)|\le C_Tl_jq(l_j)$; \textup{(b)} for $\beta\in(0,5]$, $|c_j(x)|\le C_TQ_\beta l_j^{1-\beta}$; \textup{(c)} $|c_j(x)|\le C_TQ_1$; \textup{(d)} $|c_j(x)|\le C_T(1+\ln(L/\eps_0))^5l_j$; \textup{(e)} only the order-$4$ finite piece carries $q$: $|c^{(2)}_j|\le\max(C_D,c_m)l_j$, $|l_j^2{\dm}^{(4)}_{K,j}|\le(C_4/\ln L)l_j^2(1+\ln l_j^{-1})\le Cl_j$, $|l_j^2{\dm}^{(4)}_{\mathrm{fin}}|\le C_{\mathrm{DMf}}l_jq(l_j)$. $C_T$ is free of $\eps$, $\eps_0$, $j$, $K$ and the torus.
\end{theorem}

\begin{proof}
Split $c_j=c^{(2)}_j+c^{(4)}_j$ by order. Order $2$: for $j\ge1$, ${\dm}^{(2)}_{\mathrm{fin}}+{\dm}^{(2)}_{K,j}=\Delta_{K,K}+(\Delta_{j,j}-\Delta_{K,K})=\Delta_{j,j}$ (\cref{thm:a06e:DM}(b),(e)) and $|\Delta_{j,j}|\le C_Dl_j^{-1}$; for $j=0$ it is $\dm_1$ with $|\dm_1|\le c_ml_0^{-1}$; so $|c^{(2)}_j|\le\max(C_D,c_m)l_j$. Order $4$, $(K,j)$ piece: $|{\dm}^{(4)}_{K,j}|\le C_4(K-j)$ (\cref{thm:a06e:DM}(c)) and $K-j\le\log_L(1/l_j)\le(1+\ln l_j^{-1})/\ln L$, so $l_j^2|{\dm}^{(4)}_{K,j}|\le(C_4/\ln L)l_j^2(1+\ln l_j^{-1})\le(C_4/\ln L)l_j$, as $t(1+\ln t^{-1})\le1$ on $(0,1]$. Order $4$, finite piece: $f(t)\defeq t^{-1}q(t)$ is decreasing on $(0,1]$ ($f'/f=-t^{-1}(1+5/(1+\ln t^{-1}))<0$), and $l_j\le l_K$, so \cref{thm:a06e:DMf} gives $l_j^2|{\dm}^{(4)}_{\mathrm{fin}}|\le C_{\mathrm{DMf}}l_j^2f(l_K)\le C_{\mathrm{DMf}}l_j^2f(l_j)=C_{\mathrm{DMf}}l_jq(l_j)$. Sum, with $q\ge1$: (a). (b), (c) are $q(t)\le Q_\beta t^{-\beta}$, $tq(t)\le Q_1$. (d): $l_j^2f(l_K)=l_j(l_j/l_K)(1+\ln l_K^{-1})^5\le l_j(1+\ln(L/\eps_0))^5$, since $l_j\le l_K$ and $l_K>\eps_0/L$; the other parts are $\le Cl_j$ already.
\end{proof}

\begin{corollary}[$\dm_{\mathrm{fin}}$ at fixed $\eps_0$]\label{cor:a06e:T0p}
$|\dm_{\mathrm{fin}}(x)|\le C_TL\eps_0^{-1}(1+\ln(L/\eps_0))^5$, uniformly in $\eps$, $x$ and the torus: Ba{\l}aban's ``convergent as $\eps\to0$'' \cite[(1.26)]{Ba3}, with a constant depending on $\eps_0$.
\end{corollary}

\begin{proof}
\Cref{thm:a06e:T0}(d) at $j=K$ ($c_K=l_K^2\dm_{\mathrm{fin}}$) and $l_K>\eps_0/L$.
\end{proof}

\begin{lemma}[family tolerance]\label{lem:a06e:T1}
Let $F$ be a localised renormalised family at a step $j\le K$, in any application of Proposition~1 of Steps~3(B)--3(D), and let $n_7(F)\le\bar n/2$ be the number of its non-renormalising \cite[(1.7)]{Ba3} vertices. If a bound $B_F$ on $F$ was obtained with $|c_j|\le Cl_j$ at each such vertex, then \cref{thm:a06e:T0} gives the bound $(C_T/C)^{n_7}q(l_j)^{n_7}B_F$, with constants free of $\eps_0$. Consequently: \textup{(i)} a bound $O(1)t^\kappa p(t)^Ne^{-\delta d}$ becomes $O(1)t^\kappa p(t)^Nq(t)^{\bar n/2}e^{-\delta d}$, the same polylogarithmic class ($p(t)^Nq(t)^n\le C_{N,n,\beta}t^{-\beta}$); \textup{(ii)} a conclusion ``$C_\beta t^{\kappa-\beta}$ for every $\beta>0$'' holds with the same $\kappa$ and a new $\eps_0$-free $C_\beta$; \textup{(iii)} a bound that used only $|c_j|\le C$ holds with $C_TQ_1$; \textup{(iv)} a geometric step sum $\sum_{i=k}^{K-1}t_i^aq(t_i)^n$ is $\le(1-L^a)^{-1}t_k^aq(t_k)^n$ for $a<0$ and $\le(1-L^{-a/2})^{-1}Q^n_{a/(2n)}$ for $a>0$, $n\ge1$. Every consumer of the counterterm sizes in Steps~3(B)--3(D) uses the \cite[(1.7)]{Ba3} vertex only as a positive power of $t$ (\cite[p.~421]{Ba3}; the exponents $\kappa_0>0$, $\frac12-\beta$, $\frac14-\beta$ of \cite[(3.58)]{Ba1}, \cite[(2.57), (2.116)--(2.117)]{Ba2}), so each tolerates the form \textup{(a)} of \cref{thm:a06e:T0}; where such a bound was written with an exact factor $t$ it holds with $tq(t)$.
\end{lemma}

\begin{proof}
Proposition~1 is multilinear in the coefficients $c_j$ (Step~3(E), before \cref{lem:a06e:CUT}); replace $\sup|c_j|\le Cl_j$ by $C_Tl_jq(l_j)$ at each of the $n_7$ vertices; $n_7\le\bar n/2$ since each \cite[(1.7)]{Ba3} vertex has order $\ge2$ and the family has order $\le\bar n$. (i), (ii): $p(t)=b_0(1+\ln t^{-1})^{\bar p}$, so $p^Nq^n$ is a power of $1+\ln t^{-1}$, bounded by $C_\beta t^{-\beta}$. (iii) is \cref{thm:a06e:T0}(c). (iv): for $a<0$, $q(t_i)\le q(t_k)$ for $i\ge k$ and the series is geometric; for $a>0$, $t^{a/2}q(t)^n\le Q^n_{a/(2n)}$ and $\sum_it_i^{a/2}\le(1-L^{-a/2})^{-1}$.
\end{proof}
